\subsection{Quiz: What the Gradient Does --- The Walking Robot (Page 13)}

\textbf{The Setup}

A humanoid robot is learning to walk. The reward $r(s, a)$ is the robot's forward speed.
Move forward $\to$ positive reward. Move backward $\to$ negative reward.

\textbf{Step 1 --- Collect sample episodes}

The robot tries 4 different episodes under the current policy $\pi_\theta(a_t \mid s_t)$:

\begin{center}
\begin{tabular}{lll}
Episode & What happened & Reward \\[4pt]
$\tau^1$ & Falls forwards & Positive (moved forward) \\
$\tau^2$ & Slowly stumbles forwards & Positive \\
$\tau^3$ & Steadily walks forwards & Positive \\
$\tau^4$ & Runs forwards & Positive \\
\end{tabular}
\end{center}

All four got \textbf{positive reward} because all of them moved forward.

\textbf{Step 2 --- Compute the gradient}

\[
\nabla_\theta J(\theta)
\approx
\sum_i^N \sum_{t=1}^T
\nabla_\theta \log \pi_\theta(a_t^i \mid s_t^i)
\left( \sum_{t'=t}^{T} r(s_{t'}^i, a_{t'}^i) \right)
\]

\begin{itemize}
  \item $\nabla_\theta J(\theta)$ --- the gradient of the total reward with respect to the policy
    parameters; the ``which direction to improve'' signal
  \item $\sum_i^N$ --- sum over all $N$ sampled episodes
  \item $\sum_{t=1}^T$ --- sum over every time step in an episode
  \item $\nabla_\theta \log \pi_\theta(a_t^i \mid s_t^i)$ --- how much each action's probability
    should shift
  \item $\sum_{t'=t}^T r(s_{t'}^i, a_{t'}^i)$ --- total future reward collected from time step $t$
    onwards (the ``return'')
\end{itemize}

For every action taken in every episode, multiply ``how much this action's probability should
change'' by ``how much reward followed from this point.'' Average this over all episodes to get the
gradient --- the direction to nudge the policy.

\textbf{The Surprising Answer}

Since \textbf{all four episodes got positive reward}, the gradient gives a positive push to
\textbf{all of them} --- including the falling episode ($\tau^1$). The gradient cannot tell the
difference between ``fall forward'' and ``run forward.'' It just sees ``positive reward $\to$ do
more of that.'' So the policy gets nudged toward falling forward, which is clearly not the ideal
walking behavior.

This is the core problem: \textbf{policy gradient is high-variance} --- its gradient estimates are
noisy and sensitive to the overall scale of rewards. If everything gets positive reward, everything
gets reinforced, even bad behavior.

\subsection{Introducing the Baseline (Page 14)}

\textbf{The Core Idea}

Instead of reinforcing an action whenever the reward is positive, reinforce it only when the reward
is \textbf{better than average}. This is done by subtracting a \textbf{baseline} $b$ from the
reward.

\textbf{The Modified Gradient Formula}

\[
\nabla_\theta J(\theta)
=
E_{\tau \sim p_\theta}
\left[
  \nabla_\theta \log p_\theta(\tau) \cdot (r(\tau) - b)
\right]
\]

\begin{itemize}
  \item $\nabla_\theta J(\theta)$ --- gradient of the objective (direction to improve policy)
  \item $E_{\tau \sim p_\theta}[\cdot]$ --- average over all episodes sampled from the current
    policy
  \item $\nabla_\theta \log p_\theta(\tau)$ --- how much the probability of this episode's actions
    should shift
  \item $r(\tau)$ --- total reward of episode $\tau$
  \item $b$ --- the baseline (a reference number we subtract)
  \item $(r(\tau) - b)$ --- how much better or worse this episode was compared to the baseline
\end{itemize}

Instead of asking ``did this episode get good reward?'', we now ask ``did this episode do
\textbf{better than the baseline}?'' If yes, increase the probability of those actions. If no,
decrease it. The baseline acts like a reference point --- a ``passing grade.''

\textbf{What is the baseline?}

The simplest choice is the \textbf{average reward} across all sampled episodes:

\[
b = \frac{1}{N} \sum_i r(\tau)
\]

\begin{itemize}
  \item $N$ --- total number of sampled episodes
  \item $r(\tau)$ --- reward of each episode
  \item $b$ --- their simple average
\end{itemize}

This is just computing the average score across all your robot's attempts.

\textbf{Applying it to the walking example}

Say the average reward across the 4 episodes is some value $b$. Now:
\begin{itemize}
  \item ``Running forward'' ($\tau^4$) scores well above $b$ $\to$ $(r - b) > 0$ $\to$
    get a \textbf{positive push} (do more of this)
  \item ``Falling forward'' ($\tau^1$) scores below $b$ $\to$ $(r - b) < 0$ $\to$
    get a \textbf{negative push} (do less of this)
\end{itemize}

Now the gradient correctly discourages falling and encourages running, even though both had
positive raw rewards.

\textbf{Does subtracting the baseline break anything?}

No --- and this is proved at the bottom of the slide. Subtracting a constant baseline leaves the
\textbf{expected gradient unchanged} (it is unbiased). The proof shows:

\[
E[\nabla_\theta \log p_\theta(\tau) \cdot b]
=
b\,\nabla_\theta \int p_\theta(\tau)\, d\tau
=
b\,\nabla_\theta \cdot 1
= 0
\]

The key step: $\int p_\theta(\tau)\, d\tau = 1$ because all probabilities must sum to 1. The
gradient of a constant (1) is zero. So the baseline contributes \textbf{zero} to the expected
gradient --- it does not bias the direction of learning at all. But it \textbf{reduces variance}
significantly, because individual episode rewards are now measured relative to the average instead
of in absolute terms. Smaller swings in the gradient $=$ more stable learning.

\subsection{Quiz: Sparse Rewards --- Jacket Folding (Page 15)}

\textbf{The Setup}

A robot is learning to fold a jacket. The reward has only three possible values:

\begin{center}
\begin{tabular}{ll}
Outcome & Reward \\[4pt]
Doesn't touch the jacket & 0 \\
Folded with wrinkles & 0.5 \\
Neatly folded & 1 \\
\end{tabular}
\end{center}

\textbf{Step 1 --- Sample episodes}

\begin{center}
\begin{tabular}{lll}
Episode & What happened & Reward \\[4pt]
$\tau^1$ & Doesn't touch the jacket & 0 \\
$\tau^2$ & Folds only the sleeves & $\approx 0$ \\
$\tau^3$ & Flattens but doesn't fold & 0 \\
$\tau^4$ & Folds the jacket & 1 \\
\end{tabular}
\end{center}

Three out of four episodes got zero reward. Only $\tau^4$ succeeded.

\textbf{Step 2 --- Compute the gradient with baseline}

\[
\nabla_\theta J(\theta)
\approx
\sum_i^N \sum_{t=1}^T
\nabla_\theta \log \pi_\theta(a_t^i \mid s_t^i)
\left( \sum_{t'=t}^T r(s_{t'}^i, a_{t'}^i) - b \right)
\]

\begin{itemize}
  \item $\nabla_\theta J(\theta)$ --- the gradient; tells us which direction to nudge the policy
  \item $\nabla_\theta \log \pi_\theta(a_t^i \mid s_t^i)$ --- how much each action's probability
    should shift
  \item $\sum_{t'=t}^T r(s_{t'}^i, a_{t'}^i)$ --- total future reward from step $t$ onward in
    episode $i$
  \item $b$ --- the baseline (average reward across all episodes, which is close to 0 here)
\end{itemize}

For each action taken in each episode, push that action's probability up or down based on how much
better-than-average the outcome was. Subtract the baseline $b$ so only above-average results get
reinforced.

\textbf{What happens in practice here?}

Since most episodes scored 0 and the baseline $b \approx 0$, we get $(r - b) \approx 0$ for
$\tau^1$, $\tau^2$, $\tau^3$ --- their gradient contributions are essentially \textbf{zero}. Only
$\tau^4$, which scored 1, gives a meaningful signal.

So the gradient correctly points toward folding behavior --- but the problem is it barely received
any signal at all. Three out of four episodes contributed almost nothing to learning. This is called
the \textbf{sparse reward problem}: when success is rare, most of your data is useless, the
gradient is nearly flat most of the time, and learning is painfully slow and noisy. The baseline
helped in the walking example (page 13) but did not fully fix this new problem.

\subsection{Vanilla Policy Gradient and the On-Policy Problem (Page 16)}

\textbf{The Full Algorithm}

\begin{enumerate}
  \item \textbf{Sample} episodes $\{\tau^i\}$ from the current policy $\pi_\theta(a_t \mid s_t)$
  \item \textbf{Compute} the gradient $\nabla_\theta J(\theta)$ using those episodes
  \item \textbf{Update} the parameters
  \item \textbf{Repeat} until the policy stops improving (convergence)
\end{enumerate}

\textbf{The update equation (Step 3):}

\[
\vec{\theta} = \vec{\theta} + \alpha\,\nabla_\theta J(\theta)
\]

\begin{itemize}
  \item $\vec{\theta}$ --- the policy's parameters (all the knobs controlling what actions the
    policy picks)
  \item $\alpha$ --- the learning rate; how big a step to take (too big $\to$ overshoot, too small
    $\to$ crawl)
  \item $\nabla_\theta J(\theta)$ --- the gradient; the direction that increases total reward
\end{itemize}

Take the current parameters $\vec{\theta}$, and nudge them in the direction of the gradient by a
small amount $\alpha$. Repeat. This is plain gradient ascent --- climbing uphill on the reward
surface one step at a time.

\textbf{The Big Problem: On-Policy Learning}

Vanilla policy gradient is \textbf{on-policy}, meaning the episodes used to compute the gradient
in Step 2 must have been collected by the \textbf{exact current version of the policy}. Once you
update $\theta$ in Step 3, those episodes are immediately outdated and thrown away. You can never
reuse them.

This creates a painful cycle: collect data $\to$ compute gradient $\to$ update policy $\to$
\textbf{throw all data away} $\to$ collect new data $\to$ $\ldots$

Every single gradient step requires a fresh batch of episodes. If collecting data is slow (a real
robot trying things in the physical world, or an expensive simulation), this wastes enormous
amounts of time.

The natural question is: \textbf{what if we want to reuse data collected by an older version of
the policy?} The answer is \textbf{off-policy RL using importance sampling} --- a mathematical
correction that lets you reuse old data by accounting for the fact that the old policy and the new
policy chose actions differently. This is what the next slides address.

\subsection{Off-Policy Policy Gradient (Page 17)}

\textbf{The Big Idea: Importance Sampling}

Imagine you want to know what score a student would get on a different exam, but you only have
their answers from the original exam. You can still estimate it --- but you have to \textit{adjust}
each answer's weight based on how different the two exams are. That adjustment factor is called the
\textbf{importance sampling ratio}.

Same idea here: we collected episodes using the \textit{old} policy $\theta$, but we want to
compute the gradient for a \textit{new} (updated) policy $\theta'$. We correct for this mismatch
by multiplying by the ratio of the two policies' probabilities.

\textbf{Equation 1 --- The Exact Off-Policy Gradient}

\[
\nabla_{\theta'} J(\theta')
=
\mathbb{E}_{\tau \sim p_\theta(\tau)}
\left[
  \prod_{t=1}^T
  \frac{\pi_{\theta'}(a_t \mid s_t)}{\pi_\theta(a_t \mid s_t)}
  \left( \sum_{t=1}^T \nabla_{\theta'} \log \pi_{\theta'}(a_t \mid s_t) \right)
  \left( \sum_{t'=t}^T r(s_{t'}, a_{t'}) - b \right)
\right]
\]

\begin{itemize}
  \item $\theta'$ --- the \textbf{new} policy's parameters (the one we are trying to improve)
  \item $\theta$ --- the \textbf{old} policy's parameters (the one that actually collected the
    data)
  \item $\mathbb{E}_{\tau \sim p_\theta(\tau)}[\cdot]$ --- average over episodes collected using
    the \textbf{old} policy $\theta$
  \item $\prod_{t=1}^T \frac{\pi_{\theta'}(a_t \mid s_t)}{\pi_\theta(a_t \mid s_t)}$ --- the
    \textbf{importance sampling correction}: a product of ratios, one per time step; each ratio
    asks ``how much more or less likely would the new policy have been to take this same action?''
  \item $\nabla_{\theta'} \log \pi_{\theta'}(a_t \mid s_t)$ --- how much to shift each action's
    probability under the new policy
  \item $\sum_{t'=t}^T r(s_{t'}, a_{t'}) - b$ --- future reward from step $t$ onward, minus the
    baseline
\end{itemize}

This is the REINFORCE gradient formula from before, but now with a correction factor bolted on ---
the big product in front. It uses data from the old policy, but multiplies each episode's
contribution by how differently the new policy would have behaved throughout that episode, step by
step.

\textbf{The Problem with Equation 1}

The correction factor is a \textbf{product of $T$ ratios} --- one ratio multiplied together for
every time step in the episode. If the episode is long ($T$ is large), this product becomes
unstable:

\begin{itemize}
  \item If each ratio is slightly less than 1 $\to$ the product shrinks toward \textbf{zero}
    (gradient vanishes --- no learning signal)
  \item If each ratio is slightly greater than 1 $\to$ the product explodes toward
    \textbf{infinity} (gradient blows up --- wild, unstable updates)
\end{itemize}

This is like compounding interest but in a dangerous way. Multiplying together 100 numbers that are
each 0.9 gives $0.9^{100} \approx 0.000027$ --- basically nothing.

\textbf{Equation 2 --- Switching from Trajectories to Timesteps}

\[
\nabla_{\theta'} J(\theta')
\approx
\frac{1}{N}
\sum_{i=1}^N \sum_{t=1}^T
\frac{\pi_{\theta'}(s_{i,t},\, a_{i,t})}{\pi_\theta(s_{i,t},\, a_{i,t})}
\,\nabla_{\theta'} \log \pi_{\theta'}(a_{i,t} \mid s_{i,t})
\left( \sum_{t'=t}^T r(s_{i,t'}, a_{i,t'}) - b \right)
\]

\begin{itemize}
  \item $\frac{\pi_{\theta'}(s_{i,t},\, a_{i,t})}{\pi_\theta(s_{i,t},\, a_{i,t})}$ --- the ratio
    of how likely it was to \textbf{be in state $s$ and take action $a$} under the new policy
    versus the old policy (a single ratio per timestep, not a product over all steps)
  \item $s_{i,t},\, a_{i,t}$ --- the state and action at time step $t$ in episode $i$
  \item Everything else is the same as Equation 1
\end{itemize}

Instead of one giant product across the whole trajectory, we now use a \textbf{single ratio per
time step} --- comparing the new and old policy just at that one moment. This avoids the
vanishing/exploding problem because there is no long chain of multiplications. The slide says this
is ``much less likely to explode/vanish.'' However, the ratio
$\frac{\pi_{\theta'}(s_{i,t}, a_{i,t})}{\pi_\theta(s_{i,t}, a_{i,t})}$ is the probability of
both being in state $s$ \textit{and} taking action $a$. The state part is hard to compute --- we
do not know exactly how often the new policy would visit each state, because state visits depend on
the full history of decisions. It is ``hard to measure.''

\textbf{Equation 3 --- The Final Practical Form}

\[
\nabla_{\theta'} J(\theta')
\approx
\frac{1}{N}
\sum_{i=1}^N \sum_{t=1}^T
\frac{\pi_{\theta'}(a_{i,t} \mid s_{i,t})}{\pi_\theta(a_{i,t} \mid s_{i,t})}
\,\nabla_{\theta'} \log \pi_{\theta'}(a_{i,t} \mid s_{i,t})
\left( \sum_{t'=t}^T r(s_{i,t'}, a_{i,t'}) - b \right)
\]

\begin{itemize}
  \item $\frac{\pi_{\theta'}(a_{i,t} \mid s_{i,t})}{\pi_\theta(a_{i,t} \mid s_{i,t})}$ --- the
    ratio of \textbf{just the action probabilities}: given that we are already in state $s$, how
    much more or less likely is the new policy to pick action $a$ compared to the old policy? The
    state part is dropped entirely.
  \item Everything else is the same as Equation 2
\end{itemize}

The state visitation ratio (the hard-to-measure part) is dropped, and only the action probability
ratio is kept. This is an approximation --- we are assuming the new and old policy visit states at
roughly the same frequency --- but it makes the formula practical. Both
$\pi_{\theta'}(a \mid s)$ and $\pi_\theta(a \mid s)$ are just policy outputs that can be computed
directly by plugging the state into the policy network.

\textbf{Putting It All Together}

The three equations represent a progression of trade-offs:

\begin{center}
\begin{tabular}{llll}
Equation & Correction factor & Explode/vanish? & Easy to compute? \\[4pt]
Eq.\ 1 & Product over all $T$ steps & Yes, badly & Yes \\
Eq.\ 2 & Single state-action ratio per step & No & No (state part is hard) \\
Eq.\ 3 & Single action-only ratio per step & No & \textbf{Yes} \\
\end{tabular}
\end{center}

Equation 3 is the final usable formula. It lets you reuse data from an old policy, corrects for
the mismatch using a simple per-step ratio, and is fully computable. This is the foundation for
\textbf{off-policy} algorithms that are far more data-efficient than vanilla policy gradient.

\subsection{Off-Policy Policy Gradient: Full Algorithm (Page 18)}

\textbf{The Algorithm (same three steps, different loop)}

\begin{enumerate}
  \item Sample episodes $\{\tau^i\}$ from policy $\pi_\theta(a_t \mid s_t)$
  \item Compute $\nabla_\theta J(\theta)$ using those episodes
  \item Update: $\vec{\theta} = \vec{\theta} + \alpha\,\nabla_\theta J(\theta)$
\end{enumerate}

\textbf{The update equation:}

\[
\vec{\theta} = \vec{\theta} + \alpha\,\nabla_\theta J(\theta)
\]

\begin{itemize}
  \item $\vec{\theta}$ --- the policy's parameters (all the internal settings that control
    behavior)
  \item $\alpha$ --- the learning rate; how big each step is
  \item $\nabla_\theta J(\theta)$ --- the gradient; the direction that increases reward
\end{itemize}

Nudge the parameters in the direction that increases reward, scaled by the step size $\alpha$.
This is the same update as vanilla policy gradient.

\textbf{What is actually different: the loop}

The slide shows two arrows on the algorithm:

\begin{itemize}
  \item \textbf{Red arrow} (on-policy loop): After step 3, loop back all the way to step 1 ---
    collect \textbf{new} data before doing anything else.
  \item \textbf{Orange arrow} (off-policy loop): After step 3, loop back only to step 2 ---
    recompute the gradient on the \textbf{same old data} using the importance sampling correction
    from page 17, and update again.
\end{itemize}

In vanilla policy gradient, after one update you must throw away all your data and go collect more.
In off-policy policy gradient, you can reuse the same batch and do \textbf{multiple gradient steps}
on it --- because the importance sampling ratio in the gradient formula already corrects for the
fact that the policy has changed since the data was collected.

Think of it like studying for an exam: on-policy is like solving one practice problem and then
immediately going to get a new set --- very inefficient. Off-policy is like solving the same set
of practice problems multiple times, learning more from them each pass through.

\subsection{The KL Divergence Constraint (Page 19)}

\textbf{The Problem}

When you take many gradient steps on the same batch of old data, the new policy $\theta'$ can
drift very far from the old policy $\theta$ that collected the data. But the importance sampling
correction only works well when the two policies are \textit{similar}. Once they diverge
significantly:

\begin{itemize}
  \item The old data no longer reflects the states the new policy would actually visit
  \item The gradient estimate becomes increasingly wrong
  \item You might make a bad update that seriously hurts the policy's performance --- and you might
    not even realize it until it is too late
\end{itemize}

It is like navigating using an old map. A slightly outdated map is fine. But if the roads have
completely changed, the map will lead you off a cliff.

\textbf{The Fix --- Constrain How Much the Policy Can Change}

\[
\mathbb{E}_{s \sim \pi_\theta}
\left[
  D_{\mathrm{KL}}\!\left(\pi_{\theta'}(\cdot \mid s) \;\|\; \pi_\theta(\cdot \mid s)\right)
\right]
\leq \delta
\]

\begin{itemize}
  \item $\mathbb{E}_{s \sim \pi_\theta}[\cdot]$ --- average over states that the \textbf{old}
    policy $\theta$ visited; we are measuring the difference in the situations the old policy
    actually experienced
  \item $D_{\mathrm{KL}}(\cdot \| \cdot)$ --- \textbf{KL divergence}: a mathematical ruler for
    measuring how different two probability distributions are; it equals 0 when they are identical,
    and grows larger as they diverge
  \item $\pi_{\theta'}(\cdot \mid s)$ --- the \textbf{new} policy's action choices at state $s$
    (a full probability distribution over actions)
  \item $\pi_\theta(\cdot \mid s)$ --- the \textbf{old} policy's action choices at state $s$
  \item $\delta$ --- a small threshold (your ``budget'' for how much the policy is allowed to
    change); typically a small number like 0.01
\end{itemize}

On average across all the states the old policy visited, the new policy must not be ``too
different'' from the old one. ``Too different'' is measured by KL divergence, and the constraint
says it must stay at or below the budget $\delta$. This limits how large each policy update can be,
keeping the new policy close enough to the old one that the recycled data remains trustworthy.

\textbf{Why This Matters}

This constraint is the core idea behind \textbf{Trust Region Policy Optimization (TRPO)} --- one
of the most important algorithms in modern RL. The name comes from this: you trust the old data
only within a certain region of policy space. Inside that region (where new and old policies are
close), the gradient estimate is reliable. Outside it, you cannot trust the old data anymore and
must stop.

A simpler version of this same idea later became \textbf{PPO (Proximal Policy Optimization)},
which enforces a similar limit without the full math of a trust region, and is today one of the
most widely used RL algorithms in practice --- including in training large language models like the
ones behind ChatGPT.